\section{專案介紹}

\subsection{作業目標}

本次作業的目標是使用 C++17、OpenGL 與 FreeGLUT 實作一套互動式 2D Drawing Panel。
使用者可以透過滑鼠與鍵盤繪製基本幾何圖形與文字，並利用階層式選單調整顏色、
線寬、點大小、填滿模式與字型等繪圖屬性。此外，程式也需要處理 keyboard、
mouse、window 與 motion event，並利用 Color Buffer 保存與恢復畫面內容。

本次完成的繪圖程式命名為 \textbf{Graphtoria}。
除了作業要求的基本繪圖功能外，目前也加入 Select、Undo/Redo、Grid、多視窗、
文字輸入對話框與 PPM 圖片輸出等功能。其中 Select 可以選取畫面中最上層的圖形
或文字、顯示選取範圍，並透過可復原的 Command 刪除物件。

在實作過程中，我也希望各項功能不要全部集中在 GLUT callback 與全域狀態中，
因此將視窗管理、使用者輸入、繪圖資料與畫面渲染分開處理。這些設計除了讓程式
較容易整理，也使後續加入多視窗、Undo/Redo 與物件選取等功能時，不必重新修改
整個繪圖流程。


\subsection{系統概述}

程式啟動後會先載入繪圖所需的字型資源，再建立主要繪圖視窗。主畫面中央為
Canvas，底部狀態列顯示目前使用的工具、滑鼠位置與畫布大小；在 Canvas 上按下
滑鼠右鍵則可以開啟階層式選單，切換繪圖工具或調整各種繪圖屬性。

使用者進行繪圖時，工具會先建立暫時的 draft，並隨著滑鼠移動更新內容，讓使用者
可以即時看到目前的繪製結果。當操作完成後，物件才會正式加入 Scene，並透過
Command 記錄這次修改，因此 Undo/Redo 不需要知道各種工具實際如何收集滑鼠輸入。

畫面顯示時，程式會依序繪製 Grid、Scene 中已完成的物件、尚未完成的工具預覽，
以及目前的選取範圍。滑鼠與鍵盤輸入則先由視窗轉換成程式內部的 Event，再交由
對應的 UI 元件或 Canvas 處理。整體架構與各模組之間的關係將在後續章節進一步說明。
